\clearpage
% ============================================================
\section{Objetivos}

El presente capítulo define los objetivos que orientan el desarrollo
de la investigación, estableciendo de manera clara las metas a alcanzar
y los criterios que permiten verificar su cumplimiento.

% ------------------------------------------------------------
\subsection{Objetivo General}
% ------------------------------------------------------------

Evaluar el nivel de satisfacción percibido por el nutricionista,
mediante una escala tipo Likert de cinco puntos, respecto a las
recomendaciones de planes nutricionales personalizados generadas
por un prototipo de aplicación de gestión nutricional, utilizando datos antropométricos, bioquímicos del paciente y un
catálogo de alimentos ecuatorianos.

% ------------------------------------------------------------
\subsection{Objetivos Específicos}
% ------------------------------------------------------------

\textbf{O.E.1.} Diseñar un prototipo funcional de sistema de gestión
nutricional que integre datos antropométricos y bioquímicos del
paciente con un catálogo de alimentos ecuatorianos, empleando un
modelo de lenguaje largo con base de conocimiento nutricional estructurada
para generar recomendaciones de planes personalizados.

\medskip

\textbf{O.E.2.} Medir el nivel de satisfacción del nutricionista
mediante un cuestionario validado con escala tipo Likert de cinco
puntos, aplicado tras la interacción con el prototipo.

% ------------------------------------------------------------
\subsection{Alcance del proyecto}
% ------------------------------------------------------------

El proyecto abarca el diseño y evaluación de un prototipo de sistema
de gestión nutricional orientado al contexto clínico ecuatoriano.
El alcance se delimita en función de los dos objetivos específicos
planteados y se detalla a continuación.

% ------------------------------------------------------------
\subsubsection{Alcance O.E.1: Diseño del prototipo}
% ------------------------------------------------------------

\textbf{Actividad 1.} Recopilación y estructuración del
catálogo de alimentos ecuatorianos: selección de alimentos
representativos, normalización de su composición nutricional
(proteínas, carbohidratos, lípidos, fibra y micronutrientes) y
clasificación por grupos alimentarios adaptados al contexto local.

\medskip

\textbf{Actividad 2.} Definición de la arquitectura del sistema
mediante el modelo C4 (Contexto, Contenedores, Componentes y Código)
y configuración de la base de conocimiento nutricional estructurada,
compuesta por el catálogo de alimentos ecuatorianos y los perfiles
clínicos básicos del paciente (peso, talla, edad, nivel de actividad
física). Esta base se conecta con un modelo de lenguaje de código
abierto (p.~ej., Qwen, Kimi~K, GLM o Nemotron) mediante un pipeline
de generación aumentada por recuperación~(RAG), de modo que ante una
consulta, el sistema recupera los fragmentos más relevantes y los
entrega al modelo para generar un plan alimentario justificado.

\medskip

\textbf{Actividad 3.} Desarrollo de la interfaz del sistema:
formularios de ingreso de datos antropométricos (edad, peso, talla,
IMC) y bioquímicos (glucosa, perfil lipídico), visualización de
planes nutricionales con distribución de macronutrientes y grupos
alimentarios, cálculo de gasto energético mediante fórmulas
estandarizadas, y gestión del historial clínico del paciente.

\medskip

\textbf{Actividad 4.} Documentación técnica del prototipo: descripción
de la arquitectura implementada, decisiones de diseño adoptadas durante
el desarrollo, manual de instalación del sistema y guía de uso dirigida
al nutricionista como usuario principal.

% ------------------------------------------------------------
\subsubsection{Alcance O.E.2: Medición de satisfacción}
% ------------------------------------------------------------

\textbf{Actividad 1.} Diseño y validación del cuestionario con escala
tipo Likert de cinco puntos (1~=~Muy insatisfecho; 5~=~Muy satisfecho),
definiendo dimensiones de evaluación sobre adecuación clínica de las
recomendaciones, claridad del lenguaje, facilidad de uso, confianza
en la información generada y eficiencia temporal percibida en la
elaboración del plan nutricional.

\medskip

\textbf{Actividad 2.} Definición y ejecución de casos de prueba con
perfiles genéricos de pacientes representativos de la población
ecuatoriana, caracterizados por datos como peso, talla, edad, sexo,
nivel de actividad física y enfermedades o condiciones físicas
(p.~ej., diabetes o sobrepeso); aplicando el prototipo
a cada caso y registrando las recomendaciones generadas para su
evaluación posterior por el Nutricionista.

\medskip

\textbf{Actividad 3.} Aplicación del cuestionario a/los
nutricionistas, tabulación de respuestas, cálculo del índice global
de satisfacción y análisis estadístico descriptivo para determinar
el nivel de aceptación del sistema.

% ------------------------------------------------------------
\subsection{Limitaciones}
% ------------------------------------------------------------

Las siguientes limitaciones delimitan el alcance del proyecto y
definen lo que no será desarrollado ni evaluado durante su ejecución.

No se desarrollará la integración del prototipo con sistemas de
información hospitalaria~(HIS) ni con expedientes clínicos
electrónicos~(EHR/EMR) existentes en instituciones de salud
ecuatorianas.

No se desarrollará una aplicación móvil ni una versión desplegada
en entorno de producción; el sistema se entrega exclusivamente como
prototipo funcional en entorno de desarrollo local.

No se evaluará la precisión diagnóstica clínica de las recomendaciones
generadas ni su validación médica formal; la evaluación se limita a
la satisfacción percibida por el nutricionista mediante un cuestionario
definido.

No se evaluará el desempeño del sistema frente a otros modelos de
inteligencia artificial o sistemas de recomendación nutricional
existentes.

Finalmente, el uso de modelos de lenguaje de código abierto implica
tiempos de respuesta superiores a los de modelos comerciales
accesibles vía API, dado que la inferencia se ejecuta sobre el
hardware local disponible sin optimizaciones de infraestructura. Esta diferencia en velocidad puede afectar la fluidez
de la interacción durante la evaluación, y la calidad de las
recomendaciones puede ser inferior a la de modelos propietarios
de mayor escala.

% ------------------------------------------------------------
% TABLA RESUMEN — después de limitaciones (IEEE)
% ------------------------------------------------------------

\begin{table}[h]
\centering
\caption{RESUMEN DE ALCANCE DEL PROYECTO}
\label{tab:alcance}
\begin{tabular}{|p{1.8cm}|p{3.8cm}|p{3.5cm}|p{3.5cm}|}
\hline
\textbf{Objetivo} & \textbf{Incluido} & \textbf{Excluido}
& \textbf{Producto} \\ \hline

O.G.
& Evaluación de la satisfacción del nutricionista frente a
recomendaciones del prototipo mediante escala tipo Likert.
& Validación clínica en entorno hospitalario real; puesta en
producción del sistema.
& Informe de evaluación con índice de satisfacción global
(prototipo SW). \\ \hline

O.E.1: Act.~1
& Estructuración del catálogo de alimentos ecuatorianos.
& Validación con bases de datos internacionales.
& Catálogo nutricional estructurado. \\ \hline

O.E.1: Act.~2
& Definición de la arquitectura C4 y configuración de la base de
conocimiento con pipeline RAG y modelos de código abierto.
& Despliegue en producción; integración con sistemas hospitalarios.
& Arquitectura C4 y sistema de generación de recomendaciones funcional. \\ \hline

O.E.1: Act.~3
& Interfaz clínica con distribución de macronutrientes, grupos
alimentarios, gasto energético e historial del paciente.
& Aplicación móvil; versión en producción.
& Prototipo funcional SW. \\ \hline

O.E.2: Act.~1
& Diseño y validación del cuestionario Likert con dimensión de
eficiencia temporal del nutricionista.
& Evaluación con pacientes reales.
& Instrumento de medición validado. \\ \hline

O.E.2: Act.~2
& Definición y ejecución de casos de prueba con perfiles genéricos
(peso, talla, edad, sexo, actividad física, enfermedades).
& Pruebas clínicas controladas.
& Conjunto de casos de prueba documentados. \\ \hline

O.E.2: Act.~3
& Aplicación del cuestionario y análisis estadístico descriptivo.
& Comparación con otros modelos de inteligencia artificial.
& Informe de resultados de satisfacción. \\ \hline

\end{tabular}
\end{table}
